\section{Eutectic phase behavior}
The binary eutectic phase behavior describes the chemical behavior of 
two immiscible crystals from a completely miscible solution.

On a typical eutectic phase behavior diagram, we can locate any phase situation
using temperature and composition information.

The binary eutectic phase behavior has the following properties:
\begin{itemize}
		  \item Assume a constant liquidus density. 
\end{itemize}

\subsection{Physical attributes}
We first define concentration $c_{i,j}$ as:
\begin{equation*}
		  c_{i,j} = \frac{M_{i,j}}{V_{i,j}}
\end{equation*}
where $M_{i,j}$ and $V_{i,j}$ represent mass and volume of specie $i$ in phase $j$.

Secondly, we define porosity $\phi$ as:
\begin{equation*}
		  \phi = \frac{V_{\text{porosity}}}{V_{\text{total}}}
\end{equation*}


Giving an initial distribution of $c_{i,j}$
\begin{equation*}
		  \bar{c} = \phi_f c_f + \phi_s c_s
\end{equation*}

Definition of internal energy and enthalpy:
\begin{equation*}
H = U + pV
\end{equation*}

In the phase package 

\subsection{Thermaldynamics procedures}
Assuming isochoric process during phase change.


\subsection{Some questions}


\begin{itemize}
		  \item crystals A + B ?
		  \item How to use the $(\bar{c}, \bar{E})$ pair to find porosity $\phi$?
\end{itemize}

\newpage

\section{Conservation revisit}
Revisit Conservation equations. We have the following assumptions and definitions:
\begin{itemize}
		  \item Two miscible fluids. Liqudius mixture will not change volume. Only one $V_f$.
		  \item We assume intrinsic relation between fluid densities and pressure (Temperature). $\rho_{i,f} = \rho_{i,f}(rho^0_{i,f}, P_f, E_f)$.
		  \item Fluid densities can also be determined with $\rho_{i,f} = \frac{M_{i,f}}{V_{i,f}}$ with the only one fluid volume.
		  \item Constant solid densities. $\rho_{i,i} = \frac{M_{i,i}}{V_{i,i}}$. $\rho_{i,i}$ is acquired in advance.
		  \item Porosity defined as volume fraction $\phi_f = \frac{V_f}{V}$.
		  \item Concentration average $\bar{c}_i = \phi_f c_{i,f} + (1-\phi_f) c_{i,i} = \frac{M_{i,i} + M_{i,f}}{V}$.
\end{itemize}
If assumed mix without volume change, we get $V_{i,f} = V_{f}$ here.
Then we also get concentration $c_{i,f} = \frac{M_{i,f}}{V_f} = \rho_{i,f}$.
However, it does not hold if volume changes upon mixing.

\subsection{Purpose}
We are going to try to reterive $M_1/M$ using $\bar{c}_1$ and $\rho$.
First we assume all melting situation.
$$\bar{c}_1 = \frac{M_{i}}{V} \quad\quad \rho_{1,f} = \rho_{1,f}(rho^0_{1,f}, P_f, E_f) = \frac{M_{1,f}}{V_f} \quad\quad (1-\frac{\bar{c}_1}{\rho_{1,f}}) = \frac{V}{}$$


\subsection{Question 1}
Is mix without volume change possible? 
Example: water and ethanol are two perfectly miscible fluids.
500ml ethanol into 500ml water will result in 970ml liquid.
When two liquids are mixed, the molecules of the liquids can go into the other liquid's available spaces,
which resulted in the reduction of total volume.
If mix without volume change, all molecules of one liquid fit into another liquid's available spaces.

